% sections/conclusion.tex: Section 10.
\section{Conclusion}\label{sec:conclusion}

This paper develops and applies a method for measuring household discount rates from default decisions. In an endogenous-default model the response of default to a payment due at a given date is the product of a common scale and a date-specific weight, and ratios of responses to payments at different dates identify the discount function once survival and the marginal-utility ratio are measured, without a functional form for utility or consumption data, for the households whose preferences asset-pricing data cannot reveal.

The formula is applied to a credit-bureau panel of auto loans and to the published experiments that move payments at known dates, and the two settings are consistent. Comparing two loans of the same borrower, the ratio of the far-payment response to the current-payment response, mapped through the formula decile by decile and averaged with loan weights, is an annual discount rate of $0.11$, with a set from $0.09$ to $0.13$; the rate varies non-monotonically with credit score and income, steepest near median scores and at the lowest incomes, while default rates fall more than fifty-fold across the same deciles, and the one monotone gradient is age. The experiment whose two arms move payments at two future dates in one population fixes the sign of the far-payment weight and accepts the estimate, the same across credit score and leverage; the level is the bureau's, measured at the horizon where a payment response identifies it most precisely. Ahead of the contract horizon sits a present-payment premium, nine tenths at the next payment and a third at the sixtieth, whose level payment timing identifies as the product of present bias and liquidity. The bureau's designs establish, without an assumption about who chooses which contract, that default is a strategic decision responsive to the balance at stake and to whether the lender can collect it, and that what differs across people is exposure to the default margin.
